% !TEX program = lualatex
\documentclass[aspectratio=169,professionalfonts]{beamer}
\newcommand{\slidemode}{1}  % カラー投影モード
\usepackage{beamer_template}
\usepackage{enumerate}

\newcommand{\LectureNum}{12}
\newcommand{\LectureTitle}{フーリエ変換と積分定理}
\newcommand{\CourseName}{応用数学}
\renewcommand{\TermName}{前期}

\begin{document}

\begin{frame}{第12回　フーリエ変換と積分定理}
  \begin{block}{本回の内容}
  フーリエ変換と積分定理
  \end{block}
  \vfill\hfill{\scriptsize \textcolor{gray}{第3章 フーリエ解析　§2}}
\end{frame}

\begin{frame}{定義：フーリエ変換}
  \begin{block}{}
  \medskip\textbf{フーリエ変換}　$\bigl(F(\omega) = F[f(t)]\bigr)$：
  \[
    F(\omega) = \int_{-\infty}^{\infty} f(t) e^{-i\omega t} dt
  \]
  \medskip\textbf{逆フーリエ変換}　$\bigl(f(t) = F^{-1}[F(\omega)]\bigr)$：
  \[
    f(t) = (1/2\pi) \int_{-\infty}^{\infty} F(\omega) e^{i\omega t} d\omega
  \]
  {\footnotesize \textcolor{gray}{$F(\omega)$ を $f(t)$ のフーリエ変換、 $f(t)$ を $F(\omega)$ の逆フーリエ変換という}}
  \end{block}
  \vfill\hfill{\scriptsize \textcolor{gray}{p.\,91--95}}
\end{frame}

\begin{frame}{定理：フーリエ積分定理}
  \begin{block}{}
  \textbf{条件}：
  \begin{itemize}
    \item $f(t)$ は任意の有限区間で区分的に滑らか
    \item $\int_{-\infty}^{\infty} |f(t)| dt < \infty$（絶対可積分）
  \end{itemize}
  \[
    f(t) = (1/\pi) \int_{0}^\infty [\int_{-\infty}^{\infty} f(\tau)\cos(\omega(t-\tau)) d\tau] d\omega
  \]
  \end{block}
  \vfill\hfill{\scriptsize \textcolor{gray}{p.\,91--95}}
\end{frame}

\begin{frame}{フーリエ余弦変換}
  \begin{block}{}
  \[
    Fc(\omega) = \int_{0}^\infty f(t)\cos(\omega t) dt
  \]
  \[
    f(t) = (2/\pi) \int_{0}^\infty Fc(\omega)\cos(\omega t) d\omega
  \]
  {\footnotesize $f(t)$ が偶関数の場合に対応}
  \end{block}
\end{frame}

\begin{frame}{フーリエサイン変換}
  \begin{block}{}
  \[
    Fs(\omega) = \int_{0}^\infty f(t)\sin(\omega t) dt
  \]
  \[
    f(t) = (2/\pi) \int_{0}^\infty Fs(\omega)\sin(\omega t) d\omega
  \]
  {\footnotesize $f(t)$ が奇関数の場合に対応}
  \end{block}
\end{frame}

\begin{frame}{例題1　（p.\,91）}
  \begin{exampleblock}{問題}
    $f(x) = e^{-|x|}$ のフーリエ変換を求めよ\\[2pt]
  \end{exampleblock}
\end{frame}

\begin{frame}{例題1　解答}
  \begin{exampleblock}{解答}
  \begin{enumerate}[1.]
    \item[] $F(u) = \int_{-\infty}^{\infty} e^{-|x|} e^{-iux} dx$
    \item[] $= \int_{-\infty}^{0} e^{x} e^{-iux} dx + \int_{0}^{\infty} e^{-x} e^{-iux} dx$
    \item[] $= 1/(1-iu) + 1/(1+iu)$
  \end{enumerate}
  \end{exampleblock}
  \begin{alertblock}{答}
    $F(u) = 2/(1+u^{2})$
  \end{alertblock}
  {\footnotesize \textcolor{gray}{\textit{$\lim_{x \to ±\infty} e^{(1-iu)x} = 0$ を用いる}}}
\end{frame}

\begin{frame}{例題2　（p.\,92）}
  \begin{exampleblock}{問題}
    $f(x) = \begin{cases} 1 & (|x| \leq 1) \\ 0 & (|x| > 1) \end{cases}$ のフーリエ変換を求めよ\\[2pt]
  \end{exampleblock}
\end{frame}

\begin{frame}{例題2　解答}
  \begin{exampleblock}{解答}
  \begin{enumerate}[1.]
    \item[] $F(u) = \int_{-1}^{1} e^{-iux} dx = [-e^{-iux}/(iu)]_{-1}^{1}$
    \item[] $= (e^{iu} - e^{-iu})/(iu) = 2\sin u / u$
  \end{enumerate}
  \end{exampleblock}
  \begin{alertblock}{答}
    $F(u) = 2\sin u / u$
  \end{alertblock}
\end{frame}

\begin{frame}{例題3　（p.\,95）}
  \begin{exampleblock}{問題}
    $f(x) = \begin{cases} 1-|x| & (|x| \leq 1) \\ 0 & (|x| > 1) \end{cases}$ のフーリエ余弦変換を求めよ\\[2pt]
  \end{exampleblock}
\end{frame}

\begin{frame}{例題3　解答}
  \begin{exampleblock}{解答}
  \begin{enumerate}[1.]
    \item[] $f(x)$ は偶関数なのでフーリエ余弦変換を用いる
    \item[] $F(u) = 2\int_{0}^{1} (1-x) \cos(ux) dx$
    \item[] $= 2\begin{cases} [(1-x)\sin(ux)/u]_{0}^1 + (1/u)\int_{0}^1 \sin(ux) dx \end{cases}$
    \item[] $= 2\begin{cases} 0 + (1/u)[-\cos(ux)/u]_{0}^1 \end{cases} = 2(1-\cos u)/u^{2}$
  \end{enumerate}
  \end{exampleblock}
  \begin{alertblock}{答}
    $F(u) = 2(1-\cos u)/u^{2}$
  \end{alertblock}
\end{frame}

\begin{frame}{演習問題（p.\,91--95）}
  \begin{exampleblock}{自分で解いてみよう}
  \begin{description}
    \item[問1] 次の関数のフーリエ変換を求めよ
    $(1) f(x) = { 0  (x > 0),  e^x  (x \leq 0) }$\\
    $(2) g(x) = { xe^{-x}  (x > 0),  0  (x \leq 0) }$\\
    \item[問2] 次の関数のフーリエ変換を求めよ
    $(1) f(x) = { 1  (0 \leq x \leq 1),  0  (x < 0, x > 1) }$\\
    $(2) g(x) = { x  (|x| \leq 1),  0  (|x| > 1) }$\\
    \item[問3] 問2(1)の関数に積分定理を適用して、次の等式を証明せよ
    $(1) (1/2\pi)\int_{-\infty}^{\infty} (1-e^{-iu})/iu \cdot e^{iux} du = { 1  (0<x<1),  1/2  (x=0,1),  0  (x<0, x>1) }$\\
    $(2) \int_{-\infty}^{\infty} \sin u / u  du = \pi$\\
    \item[問4] 次の関数のフーリエサイン変換を求めよ
    $f(x) = { -1  (-1 \leq x < 0),  1  (0 \leq x \leq 1),  0  (|x| > 1) }$\\
  \end{description}
  \end{exampleblock}
\end{frame}

\end{document}
